给你一个长为 $n$ 的序列 $a$，$m$ 次询问，每次查询一个区间的逆序对数。

数据范围：$1 \le n,m \le 10^5,\ 0 \le a_i \le 10^9$。

直接莫队每次转移至少是 $O(\log n)$ 的，观察可以发现我们每次转移要求的信息是「一个数在某个区间内的排名」，而这一信息关于区间可以差分，因此考虑二次离线。

我们将 $[l, r]$ 拓展到 $[l, r + 1]$，要求的是在 $[l, r]$ 区间内有多少比 $a_{r+1}$ 大的数。我们用 $f(x, r)$ 表示在 $[1, r]$ 区间内比 $a_x$ 大的数，$g(x, r)$ 表示 $[1, r]$ 区间内比 $a_x$ 小的数。则答案的变化量可以写作 $f(r + 1, r) - f(r + 1, l - 1)$。

其余指针移动同理。

对于这些式子中的 $f(x, x - 1)$ 和 $g(x, x - 1)$，我们可以使用树状数组提前 $O(n \log n)$ 预处理出来；对于其余的 $f(x, p)$ 和 $g(x, p)$ 项，我们将 $x$ 离线存在 $l$ 或 $r$ 进行批量处理。

这里有一个优化空间的技巧：我们发现处理每个询问时，离线到 $p$ 上的 $x$ 都是连续的一段，因此我们不需要把每次移动都存下，只需要存下调整的一段即可。这样我们的空间可以从总次数 $O(n \sqrt{m})$ 下降到询问数 $O(m)$。

现在我们来处理我们二次离线下来的问题：向一个集合中加入数，查询一个数在集合中的排名。莫队总共移动端点的次数是 $O(n \sqrt{m})$ 的，而数组总共只有 $O(n)$ 的长度，因此我们考虑使用 $O(\sqrt{n})$ 插入、$O(1)$ 查询的值域分块解决这个问题。

至此，我们在 $O(n \sqrt{m} + n \sqrt{n})$ 的时间复杂度和 $O(n + m)$ 的空间复杂度下解决了此题。

最后值得注意的是，我们求得的是每次的答案变化量而非答案本身，需要对其进行前缀和才能得到最终的答案。

总结就是，我们首先把询问离线到莫队上面，然后我们不直接利用莫队计算答案，而是把进行莫队过程的答案的变化量离线下来计算。